% ===========================================================================
% Section 3: Method
% ===========================================================================

This section presents \ours{} (\textbf{T}ool-\textbf{A}ugmented \textbf{R}etrieval \textbf{A}gent), our complexity-adaptive self-corrective RAG framework.
We first formalize the problem (\S\ref{subsec:formulation}), then describe the overall architecture (\S\ref{subsec:architecture}), the ReAct agent and its tools (\S\ref{subsec:agent}), the self-corrective refinement mechanism (\S\ref{subsec:refinement}), the DSPy declarative pipeline (\S\ref{subsec:dspy}), and the hybrid retrieval system (\S\ref{subsec:retrieval}).

% ---------------------------------------------------------------------------
\subsection{Problem Formulation}
\label{subsec:formulation}
% ---------------------------------------------------------------------------

Given a question $q$ and a document corpus $\mathcal{C} = \{d_1, d_2, \ldots, d_N\}$, the goal of retrieval-augmented generation is to produce an answer $a$ by first retrieving a set of relevant passages $\mathcal{P} \subset \mathcal{C}$ and then generating $a$ conditioned on both $q$ and $\mathcal{P}$:
\begin{equation}
    a = \text{Generate}(q, \mathcal{P}), \quad \mathcal{P} = \text{Retrieve}(q, \mathcal{C}).
    \label{eq:rag}
\end{equation}

In self-corrective RAG, the retrieval step is iteratively refined based on a quality assessment function $\text{Eval}(q, \mathcal{P}) \rightarrow s \in [0, 100]$ that scores the retrieved passages.
Existing approaches~\cite{yan2024crag} implement this as a fixed loop:
\begin{equation}
    \mathcal{P}^{(t+1)} = \text{Retrieve}(q^{(t+1)}, \mathcal{C}), \quad q^{(t+1)} = \text{Refine}(q^{(t)}, s^{(t)}),
    \label{eq:loop}
\end{equation}
iterating until $s^{(t)} \geq \tau$ (quality threshold) or $t \geq T$ (maximum iterations).

We replace this fixed loop with an \emph{agentic} formulation.
Let $\mathcal{T} = \{t_1, t_2, \ldots, t_K\}$ be a set of tools available to a ReAct agent~\cite{yao2023react}.
At each step $i$, the agent produces a thought $h_i$, selects a tool $t_{k_i} \in \mathcal{T}$, observes the result $o_i$, and updates its internal state:
\begin{equation}
    (h_i, t_{k_i}, o_i) = \text{ReAct}(q, \mathcal{H}_{<i}), \quad \mathcal{H}_{<i} = \{(h_j, t_{k_j}, o_j)\}_{j<i}.
    \label{eq:react}
\end{equation}

The agent terminates when it determines that the accumulated passages are sufficient, producing the final passage set $\mathcal{P}^*$ and answer $a = \text{Generate}(q, \mathcal{P}^*)$.

% ---------------------------------------------------------------------------
\subsection{Overall Architecture}
\label{subsec:architecture}
% ---------------------------------------------------------------------------

% Figure 1: Overall Architecture of AgenticRAG
\begin{figure*}[t]
\centering
\resizebox{0.95\textwidth}{!}{%
\begin{tikzpicture}[
    node distance=0.8cm and 1.2cm,
    >={Stealth[length=3mm]},
    stage/.style={rectangle, draw=black!70, fill=blue!8, rounded corners=3pt,
                  minimum width=2.6cm, minimum height=0.9cm, align=center,
                  font=\small\bfseries},
    tool/.style={rectangle, draw=black!60, fill=orange!10, rounded corners=2pt,
                 minimum width=1.8cm, minimum height=0.6cm, align=center,
                 font=\scriptsize},
    agent/.style={rectangle, draw=blue!70, fill=blue!12, rounded corners=3pt,
                  align=center, font=\small\bfseries, line width=1.2pt},
    io/.style={rectangle, draw=black!40, fill=gray!8, rounded corners=2pt,
               minimum width=2cm, minimum height=0.6cm, align=center,
               font=\small},
    arrow/.style={->, thick, black!70},
    dashedarrow/.style={->, dashed, thick, black!45},
    lbl/.style={font=\scriptsize\itshape, text=black!55},
]

% Row 1: Main pipeline flow
\node[io] (input) {Question $q$};
\node[stage, right=1.4cm of input] (preprocess) {Stage 1\\Preprocessing};
\node[lbl, above=0.05cm of preprocess] {ChainOfThought};

% Agent box — place tools first, then fit box around them
% Row 1 tools at a fixed position
\node[tool] (t1) at (9.5, 0.2) {search\_passages};
\node[tool, right=0.3cm of t1] (t2) {decompose\_query};
\node[tool, right=0.3cm of t2] (t3) {evaluate\_passages};
% Row 2 tools
\node[tool, below=0.4cm of t1] (t4) {get\_passage\_detail};
\node[tool, right=0.3cm of t4] (t5) {list\_sections};
\node[tool, right=0.3cm of t5] (t6) {get\_terminology};

% Agent box — fit around all tools with padding (extra top for subtitle)
\begin{scope}[on background layer]
\node[agent, fit=(t1)(t2)(t3)(t4)(t5)(t6),
      inner xsep=0.5cm, inner ysep=1.2cm] (agentbox) {};
\end{scope}
% Title above box, subtitle inside top area (above tools)
\node[font=\small\bfseries, text=blue!70, anchor=south] at ([yshift=0.15cm]agentbox.north) {Stage 2: ReAct Agentic Refinement};
\node[font=\scriptsize\itshape, text=black!50, anchor=south] at ([yshift=0.1cm]t1.north -| agentbox.center) {Thought $\to$ Action $\to$ Observation};

% Continue main flow
\node[stage, right=1.6cm of agentbox] (merge) {Stage 3\\Merging};
\node[stage, right=1.2cm of merge] (generate) {Stage 4\\Generation};
\node[lbl, above=0.05cm of generate] {ChainOfThought};
\node[io, right=1.4cm of generate] (output) {Answer $a$};

% Row 2: Retriever
\node[stage, below=2.0cm of agentbox, fill=green!8] (retriever) {Hybrid Retriever\\FAISS + BM25 + RRF};
\node[io, below=0.7cm of retriever] (corpus) {Document Corpus $\mathcal{C}$};

% === Arrows ===
% Main flow
\draw[arrow] (input) -- (preprocess) node[midway, above, lbl] {$q$};
\draw[arrow] (preprocess) -- (agentbox.west |- preprocess) node[midway, above, lbl] {$q', \mathcal{K}$};
\draw[arrow] (agentbox.east |- merge) -- (merge) node[midway, above, lbl] {$\mathcal{P}_{\text{agent}}$};
\draw[arrow] (merge) -- (generate) node[midway, above, lbl] {$\mathcal{P}^*$};
\draw[arrow] (generate) -- (output) node[midway, above, lbl] {$a$};

% Retriever connections
\draw[arrow] (corpus) -- (retriever);
\draw[dashedarrow] (agentbox.south) -- (retriever.north);

% Initial retrieval bypass to merge
\draw[dashedarrow] (preprocess.south) -- ++(0,-0.6) -| (merge.south)
    node[pos=0.75, right, lbl] {$\mathcal{P}_{\text{init}}$};

\end{tikzpicture}
}%
\caption{Overall architecture of \ours{}. The pipeline consists of four stages: (1)~query preprocessing, (2)~ReAct agentic refinement with six specialized tools, (3)~passage merging, and (4)~answer generation. Dashed arrows indicate the initial retrieval path and retriever connection.}
\label{fig:architecture}
\end{figure*}


% Figure 2: Pipeline Comparison — Naive vs CRAG vs Loop vs Agentic
\begin{figure*}[t]
\centering
\resizebox{0.95\textwidth}{!}{%
\begin{tikzpicture}[
    node distance=0.6cm and 0.5cm,
    >={Stealth[length=2mm]},
    box/.style={rectangle, draw=black!50, rounded corners=2pt,
                minimum width=1.6cm, minimum height=0.55cm, align=center,
                font=\scriptsize},
    pbox/.style={box, fill=blue!8},
    ebox/.style={box, fill=orange!8},
    gbox/.style={box, fill=green!8},
    abox/.style={box, fill=red!8},
    arrow/.style={->, thick, black!60},
    darrow/.style={->, dashed, thick, black!35},
    plabel/.style={font=\small\bfseries, text=black!80},
]

% === (a) Naive RAG ===
\node[plabel] (la) at (0, 0) {(a) Naive RAG};
\node[pbox, below=0.2cm of la] (a1) {Retrieve};
\node[gbox, right=0.4cm of a1] (a2) {Generate};
\draw[arrow] (a1) -- (a2);

% === (b) CRAG ===
\node[plabel] (lb) at (6, 0) {(b) CRAG};
\node[pbox, below=0.2cm of lb] (b1) {Retrieve};
\node[ebox, right=0.4cm of b1] (b2) {Evaluate\\[-1pt]{\tiny(binary)}};
\node[abox, right=0.4cm of b2] (b3) {Correct /\\[-1pt]Refine};
\node[gbox, right=0.4cm of b3] (b4) {Generate};
\draw[arrow] (b1) -- (b2);
\draw[arrow] (b2) -- (b3);
\draw[arrow] (b3) -- (b4);

% === (c) Loop Refinement ===
\node[plabel] (lc) at (0, -2.5) {(c) Loop Refinement};
\node[pbox, below=0.2cm of lc] (c1) {Preprocess};
\node[pbox, right=0.4cm of c1] (c2) {Retrieve};
\node[ebox, right=0.4cm of c2] (c3) {Evaluate\\[-1pt]{\tiny(4D)}};
\node[abox, right=0.4cm of c3] (c4) {Refine\\[-1pt]Query};
\node[gbox, right=0.4cm of c4] (c5) {Generate};
\draw[arrow] (c1) -- (c2);
\draw[arrow] (c2) -- (c3);
\draw[arrow] (c3) -- (c4) node[midway, above, font=\tiny\itshape] {low};
\draw[arrow] (c3.north) -- ++(0, 0.4) -| (c5.north) node[near start, above, font=\tiny\itshape] {pass};
\draw[darrow] (c4.south) -- ++(0, -0.4) -| (c2.south) node[near start, below, font=\tiny\itshape] {loop};

% === (d) AgenticRAG (Ours) ===
\node[plabel] (ld) at (0, -5.3) {(d) \ours{} (Ours)};
\node[pbox, below=0.2cm of ld] (d1) {Preprocess};

% Tools first, then fit box around them
\node[box, fill=purple!4, draw=purple!30, minimum width=1.1cm, font=\tiny]
    (dt1) at (4.5, -5.8) {search};
\node[box, fill=purple!4, draw=purple!30, minimum width=1.1cm, font=\tiny,
      right=0.15cm of dt1] (dt2) {decompose};
\node[box, fill=purple!4, draw=purple!30, minimum width=1.1cm, font=\tiny,
      right=0.15cm of dt2] (dt3) {evaluate};
\node[box, fill=purple!4, draw=purple!30, minimum width=1.1cm, font=\tiny,
      below=0.15cm of dt1] (dt4) {inspect};
\node[box, fill=purple!4, draw=purple!30, minimum width=1.1cm, font=\tiny,
      right=0.15cm of dt4] (dt5) {sections};
\node[box, fill=purple!4, draw=purple!30, minimum width=1.1cm, font=\tiny,
      right=0.15cm of dt5] (dt6) {terms};

% Agent box — fit around tools (background layer so tools stay visible)
\begin{scope}[on background layer]
\node[rectangle, draw=purple!40, fill=purple!6, rounded corners=2pt,
      line width=0.8pt, fit=(dt1)(dt2)(dt3)(dt4)(dt5)(dt6),
      inner xsep=0.4cm, inner ysep=0.5cm] (d2) {};
\end{scope}
\node[font=\scriptsize\bfseries, text=purple!60, anchor=south]
      at ([yshift=0.05cm]d2.north) {ReAct Agent};

\node[gbox, right=0.5cm of d2] (d3) {Generate};
\draw[arrow] (d1) -- (d2);
\draw[arrow] (d2) -- (d3);

\end{tikzpicture}
}%
\caption{Comparison of pipeline architectures. (a)~Naive RAG: single-pass retrieval. (b)~CRAG: binary evaluation with single-pass correction. (c)~Loop Refinement: fixed retrieve--evaluate--refine cycle. (d)~\ours{}: ReAct agent with six tools for adaptive refinement.}
\label{fig:pipeline_comparison}
\end{figure*}


Figure~\ref{fig:architecture} illustrates the \ours{} pipeline, and Figure~\ref{fig:pipeline_comparison} contrasts it with prior approaches. The pipeline consists of four stages:

\begin{enumerate}
    \item \textbf{Query Preprocessing.} The input question $q$ is rephrased into a standalone search query $q'$ and a set of search keywords $\mathcal{K}$ using a ChainOfThought module (DSPy \texttt{PreprocessSignature}). Optionally, a Hypothetical Document Embedding (HyDE)~\cite{gao2023hyde} query is generated for improved dense retrieval.

    \item \textbf{ReAct Agentic Refinement.} The preprocessed query is passed to a ReAct agent equipped with six tools (\S\ref{subsec:agent}). The agent iteratively searches, decomposes, evaluates, and inspects passages through autonomous reasoning. This stage replaces the fixed iterative loop in prior work.

    \item \textbf{Passage Merging.} Agent-selected passages are merged with initial retrieval results (from Stage 1) using priority-based deduplication: agent-refined passages take precedence, followed by initial results, capped at $M = 30$ passages.

    \item \textbf{Answer Generation.} A ChainOfThought generator (DSPy \texttt{QnAGenerate\-Signature}) produces the final answer conditioned on the merged passages, with adaptive response formatting based on question type.
\end{enumerate}

% ---------------------------------------------------------------------------
\subsection{ReAct Agent and Tool Design}
\label{subsec:agent}
% ---------------------------------------------------------------------------

The core of \ours{} is a ReAct agent~\cite{yao2023react} that interleaves reasoning (Thought) with tool execution (Action) and result observation (Observation).
Unlike fixed-loop refinement, the agent dynamically decides \emph{which} tools to call, \emph{how many} times, and \emph{when} to stop based on its assessment of the current retrieval state.

The agent operates through a DSPy Signature (the \texttt{AgenticRefinement\-Signature}) that specifies typed inputs and outputs.
The signature's docstring serves as a structured protocol, guiding the agent through a recommended sequence of steps while allowing deviations based on the agent's reasoning.

\subsubsection{Tool Set}
\label{subsubsec:tools}

The agent has access to six tools, each implemented as a closure over shared pipeline components (retriever, indexer, evaluator).
The first four are \emph{core tools} invoked on virtually every question; the remaining two are \emph{domain-adaptive auxiliary tools} whose utility depends on corpus characteristics (see RQ4 in \S\ref{subsec:rq4}):

\begin{enumerate}
    \item \textbf{\texttt{search\_passages}$(q, \mathcal{K})$}: Executes hybrid retrieval (FAISS + BM25 with RRF fusion) and returns passage previews (first 200 characters) to conserve context. This is the primary retrieval tool, called on average 3--4 times per question.

    \item \textbf{\texttt{decompose\_query}$(q)$}: Decomposes a complex multi-hop question into 2--4 independent sub-questions using a DSPy \texttt{DecomposeQuerySignature}. Returns a boolean \texttt{is\_multi\_hop} flag and the sub-question list.

    \item \textbf{\texttt{evaluate\_passages}$(q, \mathcal{P})$}: Assesses passage quality on four dimensions---Relevance (0--30), Coverage (0--25), Specificity (0--25), Sufficiency (0--20)---producing a total score $s \in [0, 100]$ with an action decision (\texttt{output}\slash\texttt{refine}) and refinement feedback. A mandatory fallback ensures this tool is called at least once per question.

    \item \textbf{\texttt{get\_passage\_detail}$(id)$}: Retrieves the full text of a passage by ID, enabling the agent to read complete content after identifying promising passages from search previews. This implements a progressive information disclosure strategy.

    \item \textbf{\texttt{list\_document\_sections}$(id)$}: Browses the hierarchical structure (sections, table of contents) of a source document, enabling structure-aware navigation for documents with rich organizational hierarchy.

    \item \textbf{\texttt{get\_terminology}$(term)$}: Maps user-facing terms to document-internal terminology, bridging vocabulary gaps between questions and corpus language.
\end{enumerate}

An optional \texttt{calculate} tool is available for domain-specific numerical computations (\eg, financial calculations) but is not part of the six core retrieval tools.

\subsubsection{Agent Protocol}
\label{subsubsec:protocol}

The agent's signature includes a structured protocol embedded in the docstring:

\begin{enumerate}
    \item Call \texttt{decompose\_query} to determine multi-hop structure.
    \item Call \texttt{search\_passages} with the initial query; if multi-hop, search each sub-question separately.
    \item Optionally call \texttt{get\_passage\_detail} on promising results.
    \item \textbf{Mandatory}: Call \texttt{evaluate\_passages} on the accumulated passage set.
    \item If score $< \tau$: use evaluation feedback to refine the query and search again.
    \item Call \texttt{finish} when quality is sufficient or further search is unlikely to improve results.
\end{enumerate}

The protocol is a recommendation, not a constraint---the agent may deviate based on its reasoning.
However, a mandatory evaluate fallback ensures that quality assessment is never entirely skipped: if the agent reaches \texttt{finish} without calling \texttt{evaluate\_passages}, the pipeline automatically runs the evaluation tool on the accumulated passages.

Figure~\ref{fig:react_example} shows an example agent trajectory for a 3-hop question.

% Figure 3: ReAct Agent Workflow Example
\begin{figure}[t]
\centering
\small
\begin{tikzpicture}[
    >={Stealth[length=2mm]},
    thought/.style={rectangle, draw=blue!50, fill=blue!5, rounded corners=2pt,
                    text width=5.8cm, align=left, font=\scriptsize,
                    inner sep=3pt},
    action/.style={rectangle, draw=orange!50, fill=orange!5, rounded corners=2pt,
                   text width=5.8cm, align=left, font=\scriptsize,
                   inner sep=3pt},
    obs/.style={rectangle, draw=green!50, fill=green!5, rounded corners=2pt,
                text width=5.8cm, align=left, font=\scriptsize,
                inner sep=3pt},
    stepnum/.style={circle, draw=black!40, fill=black!8, minimum size=0.4cm,
                 font=\scriptsize\bfseries},
    inbox/.style={below=0.03cm of},  % gap between boxes within a step
    stepgap/.style={below=0.35cm of}, % gap between steps
]

% Step 1 — T + A + O
\node[thought] (t1) at (3.3, 0) {\textbf{T:} Multi-hop question. Decompose first.};
\node[action, below=0.03cm of t1] (a1) {\textbf{A:} \texttt{decompose\_query}(``Who directed the film starring the actor born in Springfield?'')};
\node[obs, below=0.03cm of a1] (o1) {\textbf{O:} sub\_questions=[``actor born in Springfield?'', ``film starring actor?'', ``director?'']};
\node[stepnum] at (0, 0) {1};

% Step 2 — T + A + O
\node[thought, below=0.35cm of o1] (t2) {\textbf{T:} Search sub-questions for bridging facts.};
\node[action, below=0.03cm of t2] (a2) {\textbf{A:} \texttt{search\_passages}(``actor born in Springfield'')};
\node[obs, below=0.03cm of a2] (o2) {\textbf{O:} 5 passages. ``John Smith was born in Springfield, IL...''};
\node[stepnum] at (0 |- t2) {2};

% Step 3 — A + O
\node[action, below=0.35cm of o2] (a3) {\textbf{A:} \texttt{evaluate\_passages}(question, passages)};
\node[obs, below=0.03cm of a3] (o3) {\textbf{O:} Score 62/100, \textit{refine}. ``Missing director info.''};
\node[stepnum] at (0 |- a3) {3};

% Step 4 — A + O
\node[action, below=0.35cm of o3] (a4) {\textbf{A:} \texttt{search\_passages}(``film directed starring John Smith'')};
\node[obs, below=0.03cm of a4] (o4) {\textbf{O:} 3 passages. Score 85/100 $\to$ \texttt{finish}};
\node[stepnum] at (0 |- a4) {4};

\end{tikzpicture}
\caption{Example ReAct trajectory for a 3-hop question. T=Thought, A=Action, O=Observation. The agent decomposes, searches, evaluates, and refines based on feedback.}
\label{fig:react_example}
\end{figure}


% ---------------------------------------------------------------------------
\subsection{Self-Corrective Refinement Mechanism}
\label{subsec:refinement}
% ---------------------------------------------------------------------------

The self-corrective loop in \ours{} differs from prior work in three key aspects:

\paragraph{Adaptive Depth}
Rather than iterating a fixed number of times, the agent decides when to stop based on quality signals.
Simple questions may require only one search--evaluate cycle, while complex multi-hop questions trigger multiple decompose--search--evaluate iterations.
The maximum number of ReAct iterations is set to $I = 5$ to prevent unbounded execution.

\paragraph{Progressive Leniency}
The quality threshold decreases with each evaluation attempt to prevent over-iteration:
\begin{equation}
    \tau_{\text{eff}}^{(t)} = \tau - t \cdot \delta, \quad \delta = 5,
    \label{eq:progressive}
\end{equation}
where $\tau = 40$ is the base threshold.
This ensures the agent eventually terminates with the best available passages rather than continuing to search for marginally improved results.

\paragraph{Structured Feedback}
When the evaluation score falls below the threshold, the evaluation tool provides targeted feedback---suggested queries, keywords to add or remove---that directly informs the agent's next search action.
This feedback loop replaces generic retry logic with informed refinement.

% ---------------------------------------------------------------------------
\subsection{DSPy Declarative Pipeline}
\label{subsec:dspy}
% ---------------------------------------------------------------------------

The entire \ours{} pipeline is implemented using DSPy~\cite{khattab2024dspy}, a framework for programming LM-based systems through declarative \emph{Signatures} and composable \emph{Modules}.

\subsubsection{Signature Design}
\label{subsubsec:signatures}

Each pipeline stage is defined as a DSPy Signature with typed input and output fields:

{\small
\begin{itemize}
    \item \textbf{PreprocessSignature}: (question, history) $\rightarrow$ (rephrased\_question, keywords, topic)
    \item \textbf{EvaluationSignature}: (question, passages, retry\_count) $\rightarrow$ (scores$_{4D}$, action, feedback)
    \item \textbf{QnAGenerateSignature}: (question, passages) $\rightarrow$ (answer, footnotes)
    \item \textbf{DecomposeQuerySignature}: (question) $\rightarrow$ (is\_multi\_hop, sub\_questions)
    \item \textbf{AgenticRefinementSignature}: (question, query, keywords, $\tau$) $\rightarrow$ (passages, action)
\end{itemize}
}

The typed field descriptions serve as structured prompts, replacing hand-crafted JSON templates.
Each Signature's docstring provides task-specific instructions and scoring guidelines.

\subsubsection{Module Composition}
\label{subsubsec:modules}

Signatures are composed with DSPy modules:
\begin{itemize}
    \item \texttt{ChainOfThought}~\cite{wei2022chain}: Wraps the PreprocessSignature and QnAGenerateSignature to elicit step-by-step reasoning before outputs.
    \item \texttt{Predict}: Wraps the EvaluationSignature for direct structured output.
    \item \texttt{ReAct}: Wraps the AgenticRefinementSignature with tool integration, implementing the Thought--Action--Observation loop.
\end{itemize}

\subsubsection{Automatic Optimization}
\label{subsubsec:optimization}

DSPy enables automatic prompt optimization without modifying the pipeline code:

\begin{itemize}
    \item \textbf{BootstrapFewShot}~\cite{khattab2024dspy}: Collects successful execution traces as few-shot demonstrations. Each trace includes the full input--output chain, providing in-context examples that improve LM performance.

    \item \textbf{MIPROv2}~\cite{opsahl2024miprov2}: Generates and optimizes natural language instructions for each Signature using Bayesian surrogate models. It jointly optimizes instructions and few-shot examples across all pipeline modules.
\end{itemize}

Both optimizers use a training set of 50 examples with F1 score as the optimization metric.
The optimization cost is amortized: training is performed once, and the optimized prompts (few-shot demonstrations and/or instructions) are reused for all subsequent queries.

% ---------------------------------------------------------------------------
\subsection{Hybrid Retrieval}
\label{subsec:retrieval}
% ---------------------------------------------------------------------------

\ours{} uses a hybrid retrieval system combining dense and sparse search with Reciprocal Rank Fusion (RRF)~\cite{cormack2009rrf}:

\paragraph{Dense Retrieval}
Documents are embedded using a sentence transformer model and indexed with FAISS~\cite{johnson2019faiss} for approximate nearest neighbor search.
Given a query $q$, the dense retriever returns the top-$k$ passages by cosine similarity.

\paragraph{Sparse Retrieval}
BM25~\cite{robertson2009bm25} provides lexical matching for keyword-heavy queries, particularly effective for entity names and technical terms that may not be well-captured by dense embeddings.

\paragraph{Reciprocal Rank Fusion}
The results from both retrievers are combined using RRF:
\begin{equation}
    \text{RRF}(d) = \sum_{r \in \{dense, sparse\}} \frac{1}{\kappa + \text{rank}_r(d)},
    \label{eq:rrf}
\end{equation}
where $\kappa = 60$ is a smoothing constant and $\text{rank}_r(d)$ is the rank of document $d$ in retriever $r$'s result list.
The fused ranking is used as the final retrieval output, with the top $k = 50$ results returned per query.
